
This paper asked whether alignment methods leave detectable ''fingerprints'' across performance dimensions---whether the choice of training signal creates systematic biases beyond explicitly optimized metrics. Through proof-of-concept validation using simulated data and minimal real-model validation, we developed and tested a diagnostic framework showing that multi-dimensional signature detection is methodologically feasible.

\subsubsection{Methodology Development}

Our signature detection framework measures models across correctness (pass@k), complexity (cyclomatic complexity, AST depth), and efficiency (runtime, memory) dimensions, applying PCA-based clustering with Cohen's d effect size analysis. The proof-of-concept demonstrates that when alignment signatures are designed into simulated data, the framework reliably detects them (simulated Cohen's d=7.835, 5.$2\times$ above detection threshold). Minimal real-model validation (3 models, 10 tasks, 300 generations) confirms the framework can extract performance profiles from actual models.

\subsubsection{Empirical Findings}

The minimal real-model validation yielded measurable performance differences: execution-aligned phi-2 (2.7B) achieved 13.0\% correctness vs. 1.0\% and 0.0\% for smaller models. The execution model ranked 0.0\% on correctness (top performer), satisfying the M1 mechanism prediction. However, the $8\times$ model scale difference prevents attributing these results to alignment rather than capacity. Preference-balance pattern detection (M2) failed (53.3\% mean rank vs. $\leq 30$\% threshold), likely due to the scale confound.

\subsubsection{Critical Limitations}

Our most critical limitation is that primary validation used simulated data for methodology testing rather than full-scale real inference. While this validates that the statistical framework works, it does not validate that real models exhibit signatures at scale. The minimal real validation (10 tasks, 3 models) confirms framework functionality but lacks statistical power for robust conclusions. Model scale confounds prevent isolating alignment effects from capacity effects. Small sample sizes and single benchmark limit generalization. Temperature sensitivity remains untested.

\subsubsection{Contributions}

Despite these limitations, our methodological contributions are clear:

\begin{enumerate}
\item A framework for detecting alignment method signatures through backward inference---from model outputs across multiple dimensions to inferred implicit objectives
\item Proof-of-concept validation that the methodology successfully detects simulated signatures (simulated Cohen's d=7.835)
\item Confirmation that the framework can profile real models and identify performance patterns
\item Identification of critical requirements for full validation: matched-scale comparisons, larger samples, full-scale inference, temperature sensitivity analysis
\end{enumerate}

\subsubsection{Path Forward}

The immediate next step is clear: full-scale real-model validation (164 tasks, 30 samples, matched-scale models). This 2-4 GPU hour experiment will determine whether actual alignment methods create detectable signatures in practice. If real Cohen's d > 1.5 with matched-scale models, the methodology transitions from proposal to validated diagnostic tool. If real d < 1.5, the hypothesis requires revision.

\subsubsection{Broader Vision}

As alignment methods proliferate---DPO, RLHF, RLAIF, execution-based training, and countless variants---understanding what they implicitly optimize for becomes increasingly important, if signatures exist in practice. This work provides a proposed methodology to make invisible fingerprints visible, pending validation that the fingerprints exist in reality.

We conclude where we began: alignment methods may leave fingerprints, and we have developed a methodology to detect them. Our contribution is showing how detection could work, what it could reveal, and why it could matter---with immediate next steps clearly defined. The proof-of-concept establishes methodological feasibility; full-scale real-model validation will determine empirical reality.
